\section{Methodology}
\label{sec:methodology}

\subsection{Citation Graph}

Let $G = (V, E)$ denote a directed citation graph, where each node $v \in V$ corresponds
to a scientific article and each edge $(u, v) \in E$ indicates that article $u$ cites
article $v$. Let $A \in \{0,1\}^{|V| \times |V|}$ be the adjacency matrix.

\subsection{Operator-Based Similarity}

We define a graph operator $\mathcal{T}$ that maps node signals to node signals. Examples
include adjacency powers, normalized Laplacian operators, and diffusion kernels. Given
$\mathcal{T}$, article similarity can be defined through transformed representations,
for example
\[
S_{\text{article}} = \phi(\mathcal{T}(A)),
\]
where $\phi$ denotes a similarity-inducing transformation such as cosine similarity, inner
products after normalization, or a kernelized comparison.

\subsection{Lifting from Articles to Researchers}

Let $B \in \{0,1\}^{|V| \times |R|}$ be the article-author incidence matrix, where
$R$ denotes the set of researchers. A researcher-level similarity matrix can then be
constructed as
\[
S_{\text{researcher}} = \Psi(B^\top S_{\text{article}} B),
\]
where $\Psi$ is a normalization or debiasing operator. This formulation makes the
aggregation step explicit and separates graph propagation from author-level scaling.

\subsection{Design Considerations}

The framework highlights three methodological choices:
\begin{enumerate}
    \item the graph operator $\mathcal{T}$,
    \item the similarity transformation $\phi$,
    \item the author aggregation and normalization rule $\Psi$.
\end{enumerate}

Different selections recover classical heuristics or motivate new hybrid models.
